% Shared manuscript content for both bioRxiv and Oxford Bioinformatics versions
% This file contains only the main text content
% Metadata (title, authors, abstract, etc.) is defined in the main template files

\section{Introduction}
Reconstructing viral genomes from metagenomic sequencing data presents considerable computational challenges, particularly for viruses exhibiting extensive genetic diversity \citep{Baaijens2017-hw,Meleshko2021-gb}. This challenge is amplified in segmented viruses such as influenza, rotavirus, and bunyaviruses, where individual segments can evolve under distinct selective pressures and can undergo reassortment. While some pipelines target specific viruses and their subtypes \citep{Shepard2016-uh}, accurate genome reconstruction from samples with unknown viral strains typically requires time-consuming manual curation of contigs and reference matching \citep{de_Vries2021-po}. This makes large-scale metagenome studies or rapid outbreak responses impractical.

To address these limitations, we developed \href{https://github.com/nf-core/viralmetagenome}{nf-core/\allowbreak viralmetagenome}, a comprehensive pipeline specifically designed for untargeted viral genome reconstruction using short read sequences. The pipeline is developed using Nextflow \citep{Di-Tommaso2017-nz} within the nf-core framework \citep{Langer2025-va}, ensuring reproducibility through containerization with Docker \citep{Merkel2014-hn} and Singularity \citep{Kurtzer2017-iw}, as well as portability across computational platforms such as local desktops, high-performance clusters and cloud environments.

\section{Pipeline Description}
\href{https://github.com/nf-core/viralmetagenome}{nf-core/viralmetagenome} implements a workflow for \textit{de novo} assembly, reference matching, and iterative consensus refinement to reconstruct viral genomes without prior target knowledge. The pipeline comprises five stages: read preprocessing, metagenomic diversity assessment, contig assembly and scaffolding, iterative consensus refinement with variant analysis, and quality control (Figure~\ref{fig:pipeline-workflow}). Unless otherwise specified, the pipeline offers multiple options to accommodate established user workflows and preferences. Detailed tool descriptions are provided in Supplementary Table 1.

\begin{figure*}[htbp]
    \centering
    \includegraphics[width=\textwidth]{fig1}
    \caption{\textbf{Overview of the nf-core/viralmetagenome pipeline for untargeted viral genome reconstruction.} Default options are highlighted in bold. nf-core/viralmetagenome processes short-read data through pre-processing, metagenomic diversity assessment, \textit{de novo} assembly with multiple assemblers, scaffolding with automated reference identification and taxonomy-guided contig clustering, and iterative consensus refinement. Quality-control metrics, assembly statistics, and coverage data are summarised in MultiQC reports and overview tables for downstream analysis.}
    \label{fig:pipeline-workflow}
\end{figure*}

\subsection{Read preprocessing}
Input reads are provided via sample sheets containing sample names and paths to short-read FASTQ files. These are preprocessed using FastQC and trimmed for adapters using fastp \citep{Chen2018-tu} (default) or Trimmomatic \citep{Bolger2014-si}. Fastp offers automated adapter detection and faster processing \citep{Chen2018-tu}. UMI deduplication is implemented using HUMID \citep{Laros2024-nx} for raw reads, or UMI-tools after mapping \citep{Smith2017-nk}. Multiple sequencing runs can be merged after trimming by specifying merge group identifiers in the sample sheet. Optional filtering with bbduk \citep{Bushnell2014-qy} or PRINSEQ++ \citep{Cantu2019-vs} removes low-complexity sequences, while host reads are removed with Kraken2 \citep{Wood2019-jl}.

\subsection{Metagenomic diversity assessment}
Preprocessed reads are taxonomically classified with Kaiju \citep{Menzel2016-tz} and Kraken2 \citep{Wood2019-jl} to maximise sensitivity across diverse viral families. Results from both classifiers are visualised using Krona \citep{Ondov2011-yp}.

\subsection{\textmd{\textit{De novo}} assembly and clustering}
The assembly workflow employs a multi-assembler approach where contigs generated by all tools are pooled and processed jointly. \textit{De novo} assembly uses SPAdes \citep{Meleshko2021-gb} (RNAviral mode), MEGAHIT \citep{Li2016-sd}, and Trinity \citep{Grabherr2011-ef}, capitalizing on their distinct algorithmic strengths to maximise genome recovery across diverse viral families and read depths. Subsequently, a BLASTn \citep{Altschul1990-sy} search against the Reference Viral Database (RVDB) \citep{Chin2025-fd} identifies candidate reference genomes for each contig. All contigs and their identified reference genomes are subjected to clustering. This strategy enables contig scaffolding against the cluster centroid whether a database reference genome or a \textit{de novo} contig, thereby combining outputs from multiple assemblers into unified scaffolds.

Clustering is performed in two sequential stages. First, taxonomic pre-clustering groups contigs based on taxonomic classification using Kraken2 \citep{Wood2019-jl} and Kaiju \citep{Menzel2016-tz}, with optional filtering to focus on specific taxonomic clades for more targeted analyses. Second, nucleotide similarity clustering within taxonomic groups is performed using CD-HIT-EST \citep{Li2006-nj} (default), VSEARCH \citep{Rognes2016-ju}, MMseqs2 \citep{Steinegger2017-ci}, vRhyme \citep{Kieft2022-km}, or Mash \citep{Ondov2019-bo}. All tools are valid options, though performance may vary by dataset \citep{Zielezinski2025-vl,Steinegger2017-ci}.

Optionally, following \textit{de novo} assembly and extension, reads can be mapped to all contigs using BWA-MEM2 \citep{Vasimuddin2019-rb} (default), or Bowtie2 \citep{Langmead2019-wx}. Clusters are filtered based on the total percentage of reads mapping to all contigs within a cluster, removing assembly artefacts. For the final scaffolding step, all cluster members are mapped to the cluster representative or centroid using Minimap2 \citep{Li2018-gi}, followed by consensus calling with iVar \citep{Grubaugh2019-xd} to generate reference-assisted scaffolds.

\subsection{Consensus generation and variant calling}
The consensus module supports reference-guided analysis and scaffold refinement using user-defined reference genomes. Users can provide one or multiple references; when multiple are provided, Mash \citep{Ondov2019-bo} selects the most similar one. The selected reference then undergoes a single refinement round to generate the consensus.

In contrast, scaffold refinement employs an iterative approach, performing up to 4 cycles (default 2) to polish the consensus sequence. Both processes map reads using BWA-MEM2 \citep{Vasimuddin2019-rb} (default), or Bowtie2 \citep{Langmead2019-wx}, followed by variant calling with BCFtools \citep{Danecek2021-je} or iVar \citep{Grubaugh2019-xd} (default). Benchmarking \citep{Bassano2022-cl} indicated that BCFtools outperforms iVar in precision and recall, while iVar identified more low-frequency variants. Users can prioritise sensitivity or specificity when selecting the variant caller. Final variant calling is performed on the polished consensus to characterise single nucleotide variants.

\subsection{Consensus Quality Control}
Quality control employs CheckV \citep{Nayfach2021-wl} to assess completeness, BLASTn \citep{Altschul1990-sy} for reference similarity, and MMseqs2 \citep{Steinegger2017-ci} against the annotated database Virosaurus \citep{Gleizes2020-rq}. These analyses enable species identification, genomic segment classification, host prediction, and retrieval of additional metadata embedded within the reference databases.

The consensus refinement is evaluated through sequence alignment with MAFFT \citep{Katoh2002-ox}, which compares final consensus genomes against \textit{de novo} contigs, intermediate consensus sequences from iterative cycles, and the scaffolding reference. All tool metrics are compiled into an interactive MultiQC report \citep{Ewels2016-hs}. Additionally, key metrics are extracted from the MultiQC report and compiled into overview tables for downstream analysis.

\section{Applications}
We validated the pipeline's performance using a diverse set of real metagenomic samples containing human and plant pathogens. nf-core/viralmetagenome successfully reconstructed high-quality or near-complete consensus genomes across various viral families, including segmented viruses (Lassa virus, \allowbreak{Orthonairovirus}, Tomato spotted wilt tospovirus) and non-segmented viruses (SARS-CoV-2, West Nile virus, Potato virus Y, Youcai mosaic virus, and Monkeypox virus; Supplementary methods S1). The pipeline proved efficient, processing 28 samples in 412 CPU hours with a maximum memory footprint of 79GB on a high-performance cluster. To facilitate local deployment, we implemented a \texttt{-profile local} option, which utilizes a smaller viral RefSeq database instead of the full RVDB database for Kaiju, reducing the pipeline's peak RAM usage to approximately 18GB.

To evaluate the pipeline’s robustness in complex scenarios, we conducted two targeted experiments. First, we assessed its ability to resolve co-infections using a simulated mixture of viral genomes (Supplementary Methods S2). The pipeline successfully distinguished and reconstructed distinct co-infecting strains, provided the genetic divergence was sufficient ($\le$ 88.7\% average nucleotide identity (ANI); Supplementary Figure 1).

The second part of our benchmarking explored the critical role of the scaffolding reference in consensus genome reconstruction. We selected sequencing data from seven Lassa virus-positive human clinical samples with diverse viral loads (Supplementary Methods S3; Supplementary Figure 2). In these samples, a significantly higher sequencing depth was observed for the S segment than for the L segment (Supplementary Figure 4). To establish a benchmark, a baseline consensus genome was generated for each sample using automated reference selection from the full unclustered-RVDB (U-RVDB) reference pool. We then systematically challenged the scaffolding process by providing reference genomes (19-30 per sample; Supplementary Table 2) spanning a range of similarity to the baseline consensus (64\%--100\% ANI; Supplementary Figure 3). For each supplied reference, a new consensus genome was generated using it as the scaffolding backbone when clustered with the \textit{de novo} contigs.

Comparing these generated sequences against their respective baseline consensus highlights the utility of reference scaffolding (Figure~\ref{fig:scaffolding-impact}), particularly in low read depth scenarios when \textit{de novo} assembly was challenging. Here, we defined consensus accuracy as the Average Nucleotide Identity between each test consensus and its baseline (Supplementary Methods S3). For the high-coverage S segment (Supplementary Figure 4), \textit{de novo} assembly was sufficient to reconstruct the full genome, rendering the choice of scaffolding reference largely redundant (Figure~\ref{fig:scaffolding-impact}A). In contrast, the low-coverage L segment proved more challenging: \textit{de novo} assembly produced shorter, incomplete fragments, demonstrating that the pipeline’s ability to identify and use a closely related database reference genome was essential to bridge gaps and recover a complete segment consensus sequence.

\begin{figure*}[htbp]
    \centering
    \includegraphics[width=\textwidth]{fig2}
        \caption{\textbf{Impact of scaffolding reference similarity on consensus reconstruction accuracy.} Consensus accuracy is defined as the ANI (Eq.~1 in Supplementary Methods) between the test consensus and the baseline consensus (generated using the pipeline’s automated reference selection from the full unclustered-RVDB). Test consensus genomes were scaffolded using either a supplied database reference genome or a \textit{de novo} contig. \textbf{(A)} Boxplots comparing Consensus Accuracy (\% ANI) for the L and S segments. The plots contrast the performance of using a database reference genome as the scaffolding backbone (dark grey) against relying solely on the \textit{de novo} contig (light grey). For the low-coverage L segment, using the database reference genome sequences significantly improves accuracy, whereas for the high-coverage S segment, the \textit{de novo} route performs comparably to the database reference genome scaffolding (sequencing depths detailed in Supplementary Figure S4). \textbf{(B)} Scatter plots illustrate the relationship between reference similarity and consensus accuracy (\% ANI) for the L and S segments. Here, reference similarity is defined as the ANI between the supplied database reference genome sequence and the baseline consensus genome. The dashed line represents the line of identity ($y=x$), and points are coloured by sample source (Supplementary Figure S2). Only consensus genomes that used a database reference genome sequence as the scaffold reference are shown in panel B.}
    \label{fig:scaffolding-impact}
\end{figure*}

These findings emphasise the need to keep databases like RVDB \citep{Chin2025-fd} up-to-date. Since nf-core/viralmetagenome is primarily designed for eukaryotic viruses, users interested in bacteriophage analysis are encouraged to explore pipelines designed to target phages such as VIRify \citep{Rangel-Pineros2022-wv}, VIBRANT \citep{Kieft2020-aq}, VirSorter2 \citep{Guo2021-rf}.

\section{Conclusion}
nf-core/viralmetagenome addresses a critical need in viral genomics by providing an automated, scalable solution for untargeted viral genome reconstruction. The pipeline successfully automates the traditionally time-consuming and manual process of viral genome reconstruction from short-read metagenomic data through its integrated workflow of \textit{de novo} contig assembly, automated reference selection, clustering algorithms, and iterative refinement strategies.

Our validation results indicate that the pipeline is suitable for a wide range of eukaryotic viral families, it removes the burden of extensive manual curation, and enables consistent, high-quality genome reconstruction with reproducible results across varied computational settings.
As viral surveillance and outbreak response increasingly rely on metagenomic sequencing, automated pipelines like nf-core/viralmetagenome will be essential for the timely identification of pathogen strains. The pipeline represents a significant step forward in making viral genome reconstruction accessible to researchers without requiring extensive bioinformatics expertise, facilitating broader adoption of metagenomic approaches in viral research and public health applications.


\section*{Acknowledgments}
J.K, P.L. and L.E.K acknowledge support from the Research Foundation - Flanders (Fonds voor Wetenschappelijk Onderzoek – Vlaanderen, G005323N and G051322N, 1SH2V24N, 12X9222N). Assistance from AI language models (e.g., ChatGPT) was used to enhance the clarity of the manuscript and to support code development.

\section*{Author Contributions}
J.K. designed and implemented the pipeline, performed validation analyses, and wrote the manuscript. P.L. and L.E.K. supervised the project, and M.B. provided critical feedback. The nf-core community contributed to pipeline maintenance. All authors reviewed and approved the final manuscript.

\section*{Conflict of Interest}
The authors declare no competing interests.

